\begin{tcolorbox}[colback=blue!5!white,colframe=blue!75!black,title=Note on Syllabus]
The concepts of limit superior ($\limsup$) and limit inferior ($\liminf$) are strictly \textbf{outside the scope} of the HSC Mathematics Extension 1 and Extension 2 syllabuses. However, they provide an excellent bridge into higher mathematics (such as university-level real analysis) and serve as useful tools to "estimate a boundary" for the limits of sequences and functions.
\end{tcolorbox}

\subsubsection*{Estimating Boundaries with Limsup and Liminf}
In standard mathematics, the limit of a sequence $a_n$ as $n \to \infty$ describes the exact value the sequence settles on. However, many sequences do not settle on a single value; they might oscillate or jump around (for example, $a_n = (-1)^n$). 

Instead of asking "What is the exact limit?", we can ask "What are the ultimate upper and lower boundaries of this sequence?"
\begin{itemize}
    \item \textbf{Limit Superior ($\limsup_{n \to \infty} a_n$):} The eventual upper bound of the sequence. It describes the maximum value that the sequence continues to reach or approach as $n$ becomes arbitrarily large.
    \item \textbf{Limit Inferior ($\liminf_{n \to \infty} a_n$):} The eventual lower bound of the sequence. It describes the minimum value that the sequence continues to reach or approach as $n$ becomes arbitrarily large.
\end{itemize}

A sequence $(a_n)$ converges to a limit $L$ if and only if both its limit superior and limit inferior are equal to that same expected limit $L$. That is:
\[ \lim_{n \to \infty} a_n = L \iff \limsup_{n \to \infty} a_n = \liminf_{n \to \infty} a_n = L \]

\subsubsection*{Simple Examples}

\textbf{Example 1: A Converging Sequence}\\
Consider the sequence $a_n = \frac{1}{n}$.
The terms are $1, \frac{1}{2}, \frac{1}{3}, \frac{1}{4}, \ldots$
As $n \to \infty$, the sequence approaches $0$.
Here, the eventual upper boundary and lower boundary are the same.
\[ \limsup_{n \to \infty} \frac{1}{n} = 0 \quad \text{and} \quad \liminf_{n \to \infty} \frac{1}{n} = 0 \]
Since they are equal, the standard limit exists and $\lim_{n \to \infty} \frac{1}{n} = 0$.

\textbf{Example 2: An Oscillating Sequence}\\
Consider the sequence $a_n = (-1)^n$.
The terms are $-1, 1, -1, 1, -1, 1, \ldots$
This sequence never settles on a single converging value, so the standard limit does not exist. However, the sequence is clearly bounded.
\begin{itemize}
    \item The highest value the sequence keeps hitting as $n \to \infty$ is $1$. Thus, $\limsup_{n \to \infty} (-1)^n = 1$.
    \item The lowest value the sequence keeps hitting as $n \to \infty$ is $-1$. Thus, $\liminf_{n \to \infty} (-1)^n = -1$.
\end{itemize}

\textbf{Example 3: A Sequence with Growth and Oscillation}\\
Consider $a_n = 5 + \frac{(-1)^n}{n}$.
The sequence oscillates slightly above and below $5$, getting closer and closer to $5$.
\[ \limsup_{n \to \infty} \left(5 + \frac{(-1)^n}{n}\right) = 5 \quad \text{and} \quad \liminf_{n \to \infty} \left(5 + \frac{(-1)^n}{n}\right) = 5 \]
Both estimates boundary conditions pinch together at $5$, confirming the exact limit is $5$.
